\Askhsh{14375}{4}
\begin{ylh}{2}
    \item 3.3. Εξισώσεις 2ου Βαθμού
    \item 4.1. Ανισώσεις 1ου Βαθμού
    \item 4.2. Ανισώσεις 2ου Βαθμού
    \item 5.3. Γεωμετρική πρόοδος
\end{ylh}
Δίνεται η συνάρτηση $f(x) = x^{2} - \mu x - 2,\ \ \mu\mathbb{\in R}$.
\begin{alist}
\item Να αποδείξετε ότι η εξίσωση $f(x) = 0$ έχει δύο ρίζες πραγματικές άνισες για κάθε $\mu\mathbb{\in R}$.
\monades{6}
\item Να βρείτε τις τιμές του$\ \ \mu\mathbb{\in R}$ για τις οποίες οι αριθμοί $x = - 2$ και $x = 3$ βρίσκονται εκτός του διαστήματος των ριζών της εξίσωσης $f(x) = 0$ ενώ ο $x = 1$ βρίσκεται εντός του διαστήματος των ριζών της εξίσωσης $f(x) = 0$.\monades{6}
\item Αν επιπλέον οι τιμές $f( - 2),\ \ f(1),\ \ f(3)$ με τη σειρά που δίνονται αποτελούν διαδοχικούς όρους γεωμετρικής προόδου τότε:
\begin{rlist}
\item
  Να βρείτε τις τιμές του $\mu$.
\monades{7}
\item
  Για $\mu = \dfrac{13}{7}$ να βρείτε το λόγο της παραπάνω γεωμετρικής προόδου.
\monades{6}
\end{rlist}
\end{alist}

